% V2原书思考题13.1--13.11，PDF物理页60；图13.15同轴线待精确重绘。
\section*{第13章\quad 电势：思考题}
\paragraph{思考题 13.1}
判断并举例说明：(1) 场强相等区域电势必相等；(2) 场强为零处电势必为零；(3) 电势为零处场强必为零；(4) 场强大的位置电势必较高。

\paragraph{思考题 13.2}
在一个孤立正点电荷与负点电荷的场中，用电势定义说明电势符号及沿远离点电荷方向的变化。

\paragraph{思考题 13.3}
选择从电势定义出发的推导说明静电偶极子垂直平分面上的电势为什么为零。

\paragraph{思考题 13.4}
根据静电场为电势的负梯度这一关系，说明为什么静电场线不能自身形成闭合回路。

\paragraph{思考题 13.5}
若某三维开区域内电势处处为常数，电场为什么为零？若只有某曲面上电势恒定，曲面上的电场是否一定为零？

\paragraph{思考题 13.6}
处于同一条静电电场线上的两个不同点，其电势是否可能相同？说明条件。

\paragraph{思考题 13.7}
放在较高静电势位置的电荷是否一定有较高的电势能？

\paragraph{思考题 13.8}
地球表面电场方向大致竖直向下，地面之上的电势随高度增加而怎样变化？

\paragraph{思考题 13.9}
若已知某点电场强度E，能否直接求出该处绝对电势？计算电势还需要哪些信息？

\paragraph{思考题 13.10}
小鸟停在约30 kV高压输电线上，为什么可能不被电击？它在何种情况下又会危险？

\paragraph{思考题 13.11}
原图13.15近轴同轴输电线开路端，说明内部导体和外导体之间的电势分布及电场线怎样在电缆末端向外弯曲。
